% !TeX root = ../papers/nc-p3.tex
\section{単元曲面全体での五つの中心の独立性と六次中心商}
\label{sec:rank-fifteen-bivariate-sixth-cut}
\begin{theorem}
前節の開集合の各表示において、任意の独立な四つの四次中心を課した商は、
四始点の四次と五次でそれぞれ31,40次元である。
さらに任意の幾何単元対 $X\ne1$, $Y\ne0$ に対し、六次が全始点で44次元となる
四中心と表示が取れる。$Y\ne(1-X)^2$ かつ $X^2+X+2\ne0$ なら、前節の
文字0の五自由列のうち第1,2,4,5列を、同じ表示の四始点で使える。
\end{theorem}
\begin{proof}
$F=\mathbb F_3(K,H)$ とし、前節の五本の四次中心を同じ規格化語基底で書く。
\path{computations/verify_rank_fifteen_bivariate_sixth_cut.py} は、各始点でこれらを四次経路へ
評価する行列、および各矢を右に掛けて五次経路へ評価する行列を作る。
全部で $4(1+4)=20$ 個あり、五行の階数を調べればよい。
先に開集合の十因子しか持たない主元で消去する。残る行列の最大小行列式の分子から十因子を除き、
そのイデアルに1が入るかをGröbner基底で検査する。必要なときは開集合の積 $C$ に対して
$1-zC$ も加えて局所化する。どの行列も全点で階数5である。
従って任意の四次元部分空間が各始点で四本の独立な関係を与える。
五次では四つの矢の積が別々の文字にあり、各文字で四本が独立である。
中心性により左からの積もこの関係に含まれ、$35-4=31$, $56-16=40$ を得る。

次に定理の四本を固定し、先頭語から一矢を末尾に足す形で中心商の語基底を構成する。
二次関係と四次での中心切断に、前節と同じ $P=W^2$ の偶数指数補正を使う。
例えば末尾の二次関係の係数を $W^e f$、その語を $ab$、前置する語を $u$ とする。
$ua$ の射影の基底語を $s$ としたとき、列 $sb$ の $W$ 指数は
$e-\ell_p(ua)+\ell_p(s)-\ell_p(sb)$ である。行の共通偶奇を除いて $F$ 上で消去する。
四始点の六次文字別次元はいずれも $(12,12,12,8)$ となる。

主元が消える軌道を二変数で正確に制限する。$K=1/v$, $H=c(1-v^2)/v$,
$X=c^2$, $B=(v^3-v)^2$, $Y=X^3B$ と置く。四始点の除外因子を引き戻し、
もとの開集合ですでに単元な因子を除いて最小公倍式 $G(v,c)$ を取る。
五次まで、六次までの $(v,c)$ 次数はそれぞれ $(7,7),(55,44)$ である。
$\mathbb F_3[c,v]$ は $\mathbb F_3[X,B]$ 上で $c^i v^j$ ($i<2$, $j<6$) を基底に持つ。
実際 $c^2=X$, $v^6=B+2v^4-v^2$ の二つのモニックな式で順に割れる。
$XB\ne0$ では十二表示は相異なる。従って全表示が $G=0$ となる必要十分条件は
この基底での全係数が同時に零となることである。

非零係数の対から $B$ を消去した終結式の共通因子で、必要な $X$ の条件を得る。
検証コードで選ぶ対のモニック最大公約元は、五次まででは1、六次まででは
\[
                     X^{136}(X-1)^9(X^2+X+2)
\]
である。これは二変数曲面上の孤立した不良点も検出する必要条件である。
仮定の下で残る二次因子の根は第\ref{sec:rank-fifteen-quadratic-verticals}節で
各水準に十三次までの証明書がある。分岐 $Y=(1-X)^2$ も
第\ref{sec:rank-fifteen-prime-parabolic-pencils}節で全水準を認証した。以上で六次までの主張が従う。
\end{proof}
